\section{本章小结}
\label{sec:cexp-summary}

本章从一个现有公式算不出的问题出发：只知道$\lvert X\rvert$时，$X^{+}$的平均是多少。每个条件事件的概率都是零，没有概率表也没有联合密度，组内平均无从相除。解决的办法是换一种描述：不再除以组的概率，而是要求条件平均只用已知信息，并在每个已知事件上保留原变量的累计期望。有限分组时，这两条要求恰好刻画了组内平均；一般信息下，它们就是条件期望的定义，“只知绝对值”的答案$\lvert X\rvert/2$由它直接验证。

这个定义总有解，而且解只确定到零概率集。存在性来自Radon--Nikodym定理：把$A\mapsto\mathbb{E}[X\mathbf{1}_{A}]$看作已知事件上的测度，条件期望就是它相对于$\mathbb{P}$、并且只用已知信息的密度。唯一性只用到积分定义在足够多的事件上成立。给定观测$Y$时，条件期望可以写成$m(Y)$，但函数$m$只在$Y$的分布下几乎处处确定；单个零概率观测值上的条件平均不由定义决定，要在第~\ref{chap:conditional-distributions}章借助额外的结构处理。

所有运算规则都从积分定义核验，并且都按几乎必然理解。表~\ref{tab:cexp-summary}按传感器表汇总了这些规则、它们需要的条件和例中的数值。

\begin{BookInlineTable}
  \centering
  \caption{条件期望的规则（传感器表，$m(X)=\mathbb{E}[Y\mid X]$）}
  \label{tab:cexp-summary}
  \begin{tabular}{lll}
    \toprule
    规则 & 条件 & 例中的数值 \\
    \midrule
    线性与保序 & 各变量可积 & $\mathbb{E}[2Y-1\mid X]=2m(X)-1$ \\
    提取已知量 & $H$已知，$X$与$HX$可积 & $\mathbb{E}[XY\mid X=2]=\tfrac{10}{3}$ \\
    独立的信息 & $X$与$\mathcal{G}$独立 & $\mathbb{E}[\operatorname{sgn}X\mid\lvert X\rvert]=0$ \\
    全期望公式 & $X$可积 & $0.3\cdot\tfrac13+0.4+0.3\cdot\tfrac53=1$ \\
    塔式法则 & $\mathcal{H}\subseteq\mathcal{G}$ & $\{X\leq1\}$上为$\tfrac57$ \\
    条件Jensen & $\varphi$凸，$X$与$\varphi(X)$可积 & $1\leq\mathbb{E}[Y^{2}\mid X=1]=\tfrac32$ \\
    $L^{p}$收缩 & $X\in L^{p}$，$1\leq p<\infty$ & 范数不增，逐点可增 \\
    最小均方预测 & $X\in L^{2}$ & 最小误差$\tfrac13<\tfrac{43}{105}$ \\
    全方差公式 & $X\in L^{2}$ & $0.6=\tfrac13+\tfrac{4}{15}$ \\
    \bottomrule
  \end{tabular}
\end{BookInlineTable}

规则的边界同样重要。只有已知的乘子能提到条件期望之外。条件平均为常数不能推出独立，它只约束一个平均值，不约束整个条件规律。迭代条件期望只有在两份信息有包含关系时才能化简，结果总是较粗的那一层；互不包含时两种顺序的结果不同，也不等于共同信息上的平均。凸性决定Jensen不等式的方向，证明中只能使用可数条支撑线，以便合并例外的零概率集。

条件期望的条件有两个层次。积分定义、存在唯一性、运算规则、迭代和Jensen不等式只要求可积；最小均方预测、正交投影、条件方差和全方差公式要求平方可积。在平方可积的情形，条件期望是$L^{2}$中到已知信息子空间的正交投影，塔式法则是逐次投影，全方差公式是勾股定理。信息越多，最小均方误差与平均条件方差越小，但这只是平均意义下的结论，个别样本点上的误差和个别组的条件方差都可能变大。

本章只为每个可积变量求条件平均。对每个事件$\{X\in B\}$的示性函数求条件期望，得到一族条件概率；把这一族组织成对每个观测值都是概率分布的对象，就是第~\ref{chap:conditional-distributions}章的条件分布，那里再说明条件期望怎样写成对条件分布的积分。信息随时间逐步增加、一列条件期望怎样演化，则属于随机过程的研究。
